\chapter[Band 21: Folgen]{Band 21 auf einen Blick}
\noindent\textbf{Folgen}\par\medskip
\noindent\emph{Lesart.}
Eine Familie ist eine Funktion mit ausdrücklich angegebener Indexmenge.
Bei reellen Folgen sind \(a,b\colon\mathbb N\to\mathbb R\) und
\(L,M,c\in\mathbb R\); \(\varepsilon>0\) bezeichnet eine positive reelle
Zahl. Eine Teilfolgenabbildung \(\varphi\colon\mathbb N\to\mathbb N\)
ist stets strikt wachsend. Die letzten vier Zeilen verwenden die
eigenständigen Bezeichnungen der Mogiljanskaja-Familien.

\begin{BandUebersicht}
\textbf{Indizierte Familie und Wertemenge}
Für \(a\colon I\to M\):
\UebFormel{\begin{aligned}
(a_i)_{i\in I}&\coloneqq a,\qquad a_i\coloneqq a(i),\\
\operatorname{Bild}(a)&\coloneqq a[I].
\end{aligned}}
\UebQuellen{\FormulaRefAuto{IndexedFamilyDef}[def];
\FormulaRefAuto{IndexedFamilyImageDef}[def]}
&
\textit{Elementkriterium und surjektive Familien}
\UebFormel{\begin{aligned}
&i\in I\vdash a_i\in M,\\
&x\in\operatorname{Bild}(a)
\leftrightarrow\exists i\in I\;(x=a_i),\\
&\operatorname{Bild}(a)\subseteq M,\\
&a\colon I\sur M\leftrightarrow\operatorname{Bild}(a)=M.
\end{aligned}}
\UebQuellen{\FormulaRefAuto{IndexedFamilyCoordinateTyping}[thm];
\FormulaRefAuto{IndexedFamilyImageMembership}[thm];
\FormulaRefAuto{IndexedFamilyImageSubsetCodomain}[thm];
\FormulaRefAuto{IndexedFamilySurjectiveIffImage}[thm]}
\\ \addlinespace[.8ex]

\textbf{Einschränkung und Umindizierung}
Für \(J\subseteq I\), bzw. \(\psi\colon J\to I\):
\UebFormel{\begin{aligned}
(a_i)_{i\in J}&\coloneqq a\restriction_J,\\
(a_{\psi(j)})_{j\in J}&\coloneqq a\circ\psi.
\end{aligned}}
\UebQuellen{\FormulaRefAuto{IndexedFamilyRestrictionDef}[def];
\FormulaRefAuto{IndexedFamilyReindexingDef}[def]}
&
\textit{Werte bleiben über die Indizes kontrollierbar}
\UebFormel{\begin{aligned}
&(a\restriction_J)\colon J\to M,\\
&j\in J\vdash(a\restriction_J)_j=a_j,\\
&\operatorname{Bild}(a\circ\psi)=a[\psi[J]].
\end{aligned}}
Familien mit derselben Indexmenge sind genau dann gleich, wenn
alle entsprechenden Glieder übereinstimmen.
\UebQuellen{\FormulaRefAuto{IndexedFamilyRestrictionFacts}[thm];
\FormulaRefAuto{IndexedFamilyReindexingImage}[thm];
\FormulaRefAuto{IndexedFamilyExtensionality}[thm]}
\\ \addlinespace[.8ex]

\textbf{Endliche Indexkonventionen}
Für \(n\in\mathbb N\):
\UebFormel{\begin{aligned}
&\NatSeg n=\{i\in\mathbb N\mid i<n\},\\
&\PosNatSeg n\coloneqq
\{i\in\mathbb N\mid1\leq i\leq n\},\\
&\lambda_n\colon\NatSeg n\to\PosNatSeg n,\\
&\lambda_n(i)\coloneqq i+1.
\end{aligned}}
Endliche Folgen haben eine dieser beiden Mengen als Definitionsbereich.
\UebQuellen{\FormulaRefAuto{ZeroBasedFiniteSequenceDef}[def];
\FormulaRefAuto{PositiveNaturalSegmentDef}[def];
\FormulaRefAuto{PositiveFiniteSequenceDef}[def];
\FormulaRefAuto{FiniteSequenceIndexShiftDef}[def]}
&
\textit{Bijektive Indexverschiebung}
\UebFormel{\lambda_n\colon\NatSeg n\bij\PosNatSeg n.}
Für \(a\colon\NatSeg n\to M\) ist
\(a\circ\lambda_n^{-1}\colon\PosNatSeg n\to M\) die entsprechende
positiv indizierte Folge; zurück führt Komposition mit \(\lambda_n\).
\UebQuellen{\FormulaRefAuto{FiniteSequenceIndexShiftBijective}[thm];
\FormulaRefAuto{FiniteSequenceConventionConversion}[thm]}
\\ \addlinespace[.8ex]

\textbf{Unendliche Folgen, Beispiele und Rekursionsdaten}
\UebFormel{\begin{aligned}
&(a_n)_{n\in\mathbb N}\coloneqq a
\quad(a\colon\mathbb N\to M),\\
&(\operatorname{const}_{I,c})_i\coloneqq c,\\
&(\iota_{\mathbb N})_n\coloneqq n.
\end{aligned}}
Für eine Schrittvorschrift sind \(x_0\in M\) und \(F\colon M\to M\)
gegeben.
\UebQuellen{\FormulaRefAuto{InfiniteSequenceDef}[def];
\FormulaRefAuto{ConstantIndexedFamilyDef}[def];
\FormulaRefAuto{NaturalEnumerationSequenceDef}[def]}
&
\textit{Eindeutige Rekursion und verschachtelte Mengenfolgen}
\UebFormel{\begin{aligned}
&\exists!a\colon\mathbb N\to M:\\
&\qquad a_0=x_0,\quad
\forall n\in\mathbb N\;a_{n+1}=F(a_n).
\end{aligned}}
Für eine Mengenfolge \(S\) gilt:
\UebFormel{\begin{aligned}
&\forall k\in\mathbb N\;(S_k\subseteq S_{k+1}),\quad m\leq n\\
&\qquad\vdash S_m\subseteq S_n.
\end{aligned}}
\UebQuellen{\FormulaRefAuto{RecursiveSequenceExistenceUnique}[thm];
\FormulaRefAuto{SubsetSequenceNesting}[thm]}
\\ \addlinespace[.8ex]

\textbf{Restfolge und Teilfolge}
\UebFormel{\begin{aligned}
a^{\langle r\rangle}_n&\coloneqq a_{n+r}
\quad(r\in\mathbb N),\\
(a_{\varphi(n)})_{n\in\mathbb N}&\coloneqq a\circ\varphi.
\end{aligned}}
Die strikte Monotonie von \(\varphi\) ist Teil der
Teilfolgendefinition.
\UebQuellen{\FormulaRefAuto{SequenceTailDef}[def];
\FormulaRefAuto{SubsequenceDef}[def]}
&
\textit{Grenzwert bleibt erhalten}
\UebFormel{\begin{aligned}
&a_n\to L\leftrightarrow a^{\langle r\rangle}_n\to L,\\
&a_n\to L\vdash a_{\varphi(n)}\to L.
\end{aligned}}
Auch endlich viele Änderungen der Glieder verändern den Grenzwert nicht.
\UebQuellen{\FormulaRefAuto{SequenceTailSameLimit}[thm];
\FormulaRefAuto{SubsequenceOfConvergentSequence}[thm];
\FormulaRefAuto{FiniteModificationPreservesLimit}[thm]}
\\ \addlinespace[.8ex]

\textbf{Konvergenz und Beschränktheit}
\UebFormel{\begin{aligned}
&a_n\to L\coloneqq\\
&\quad\forall\varepsilon>0\;\exists N\in\mathbb N\\
&\qquad\forall n\geq N\;|a_n-L|<\varepsilon,\\
&\operatorname{Beschr}(a)\coloneqq\\
&\quad\exists C\geq0\;\forall n\in\mathbb N\;|a_n|\leq C.
\end{aligned}}
\UebQuellen{\FormulaRefAuto{RealSequenceConvergenceDef}[def];
\FormulaRefAuto{BoundedRealSequenceDef}[def]}
&
\textit{Eindeutigkeit, Beschränktheit und Ordnung}
\UebFormel{\begin{aligned}
&a_n\to L,\ a_n\to M\vdash L=M,\\
&a_n\to L\vdash\operatorname{Beschr}(a),\\
&a_n\to L,\ b_n\to M,\ \forall n\;a_n\leq b_n\\
&\qquad\vdash L\leq M.
\end{aligned}}
\UebQuellen{\FormulaRefAuto{RealSequenceLimitUnique}[thm];
\FormulaRefAuto{ConvergentRealSequenceBounded}[thm];
\FormulaRefAuto{RealSequenceLimitPreservesOrder}[thm]}
\\ \addlinespace[.8ex]

\textbf{Nullfolge}
\UebFormel{a\text{ ist Nullfolge}\coloneqq a_n\to0.}
Mit punktweiser Subtraktion ist \((a-L)_n\coloneqq a_n-L\).
\UebQuellen{\FormulaRefAuto{NullSequenceDef}[def];
\FormulaRefAuto{PointwiseRealSequenceOperationsDef}[def]}
&
\textit{Abschätzung und Einschließung}
\UebFormel{\begin{aligned}
&a_n\to L\leftrightarrow a_n-L\to0,\\
&b_n\to0,\ \forall n\;|a_n|\leq b_n\vdash a_n\to0,\\
&u_n\to L,\ v_n\to L,\ \forall n\;u_n\leq a_n\leq v_n\\
&\qquad\vdash a_n\to L.
\end{aligned}}
Hier sind auch \(u,v\) reelle Folgen.
\UebQuellen{\FormulaRefAuto{ConvergenceIffDifferenceNullSequence}[thm];
\FormulaRefAuto{NullSequenceMajorantCriterion}[thm];
\FormulaRefAuto{RealSequenceSqueezeTheorem}[thm]}
\\ \addlinespace[.8ex]

\textbf{Monotone Folgen}
\UebFormel{\begin{aligned}
\operatorname{Mon}_{\uparrow}(a)&\coloneqq
\forall n\in\mathbb N\;(a_n\leq a_{n+1}),\\
\operatorname{Mon}_{\downarrow}(a)&\coloneqq
\forall n\in\mathbb N\;(a_{n+1}\leq a_n).
\end{aligned}}
\UebQuellen{\FormulaRefAuto{MonotoneRealSequenceDef}[def]}
&
\textit{Monotone Konvergenz}
\UebFormel{\begin{aligned}
&\operatorname{Mon}_{\uparrow}(a),\
\exists C\;\forall n\;(a_n\leq C)\\
&\qquad\vdash a_n\to\sup_{\mathbb R,\RealLe}a[\mathbb N],\\
&\operatorname{Mon}_{\downarrow}(a),\
\exists C\;\forall n\;(C\leq a_n)\\
&\qquad\vdash a_n\to\inf_{\mathbb R,\RealLe}a[\mathbb N].
\end{aligned}}
\UebQuellen{\FormulaRefAuto{MonotoneRealSequenceConvergence}[thm]}
\\ \addlinespace[.8ex]

\textbf{Punktweise Rechenoperationen}
\UebFormel{\begin{aligned}
(a+b)_n&\coloneqq a_n+b_n,\\
(ab)_n&\coloneqq a_nb_n,\qquad(ca)_n\coloneqq ca_n,\\
(a/b)_n&\coloneqq a_n/b_n
\quad(\forall n\;b_n\neq0).
\end{aligned}}
\UebQuellen{\FormulaRefAuto{PointwiseRealSequenceOperationsDef}[def];
\FormulaRefAuto{PointwiseRealSequenceReciprocalQuotientDef}[def]}
&
\textit{Grenzwertregeln}
Unter \(a_n\to L,\ b_n\to M\):
\UebFormel{\begin{aligned}
&a_n+b_n\to L+M,\qquad ca_n\to cL,\\
&a_nb_n\to LM,\qquad |a_n|\to|L|,\\
&M\neq0,\ \forall n\;b_n\neq0
\vdash a_n/b_n\to L/M.
\end{aligned}}
\UebQuellen{\FormulaRefAuto{RealSequenceLinearLimitRules}[thm];
\FormulaRefAuto{RealSequenceProductLimitRule}[thm];
\FormulaRefAuto{AbsoluteValuePreservesSequenceLimits}[thm];
\FormulaRefAuto{RealSequenceQuotientLimitRule}[thm]}
\\ \addlinespace[.8ex]

\textbf{Cauchy-Folge}
\UebFormel{\begin{aligned}
&\operatorname{Cauchy}(a)\coloneqq\\
&\quad\forall\varepsilon>0\;\exists N\in\mathbb N\\
&\qquad\forall m,n\geq N\;|a_m-a_n|<\varepsilon,\\
&\operatorname{Konv}(a)\coloneqq
\exists L\in\mathbb R\;(a_n\to L).
\end{aligned}}
\UebQuellen{\FormulaRefAuto{RealCauchySequenceDef}[def];
\FormulaRefAuto{RealSequenceConvergentDivergentDef}[def]}
&
\textit{Cauchy-Kriterium der reellen Zahlen}
\UebFormel{\begin{aligned}
&\operatorname{Konv}(a)\leftrightarrow\operatorname{Cauchy}(a),\\
&\operatorname{Cauchy}(a)\vdash\operatorname{Beschr}(a).
\end{aligned}}
Die Existenz eines Grenzwerts ist hier eine Folgerung aus der
Vollständigkeit der reellen Ordnung.
\UebQuellen{\FormulaRefAuto{RealSequenceCauchyCriterion}[thm];
\FormulaRefAuto{RealCauchySequenceBounded}[thm]}
\\ \addlinespace[.8ex]

\textbf{Mogiljanskaja: Grundschichten und Familien}
Von hier an bezeichnen \(a,b\) die folgenden Mengenfamilien:
\UebFormel{\begin{aligned}
L&\coloneqq\{0\}\times\mathbb N,\qquad a_i\coloneqq(0,i),\\
D&\coloneqq\{1\}\times(\mathbb N\times\mathbb N),\\
b_{ij}&\coloneqq(1,(i,j)).
\end{aligned}}
\UebQuellen{\FormulaRefAuto{MogiljanskajaLayerLDef}[def];
\FormulaRefAuto{MogiljanskajaLayerDDef}[def];
\FormulaRefAuto{MogiljanskajaAElementsDef}[def];
\FormulaRefAuto{MogiljanskajaBElementsDef}[def]}
&
\textit{Indizes, Bildmengen und Trennung}
\UebFormel{\begin{aligned}
&a_i=a_j\vdash i=j,\\
&b_{ij}=b_{mn}\vdash i=m,\ j=n,\\
&a[\mathbb N]=L,\qquad
b[\mathbb N\times\mathbb N]=D,\\
&L\cap D=\varnothing.
\end{aligned}}
\UebQuellen{\FormulaRefAuto{MogiljanskajaAIndexUnique}[thm];
\FormulaRefAuto{MogiljanskajaBIndexUnique}[thm];
\FormulaRefAuto{MogiljanskajaLayerLImage}[thm];
\FormulaRefAuto{MogiljanskajaLayerDImage}[thm];
\FormulaRefAuto{MogiljanskajaLayersDisjoint}[thm]}
\\ \addlinespace[.8ex]

\textbf{Marker und Reserve}
\UebFormel{\begin{aligned}
&c_0\coloneqq b_{00},\quad c_1\coloneqq b_{11},
\quad c_2\coloneqq b_{10},\\
&U\coloneqq\{b_{n+2,j}\mid n,j\in\mathbb N\},\\
&V\coloneqq U\cup\{c_2\},\quad P\coloneqq\{c_0,c_1\},\\
&D'\coloneqq P\cup V.
\end{aligned}}
\UebQuellen{\FormulaRefAuto{MogiljanskajaDPrimeDef}[def]}
&
\textit{Ausgesparte Punkte und disjunkter Kern}
\UebFormel{\begin{aligned}
&c_0,c_1,c_2\in D,\qquad c_0,c_1,c_2\notin U,\\
&c_0,c_1,c_2\text{ sind paarweise verschieden},\\
&D'\subseteq D,\qquad b_{01}\notin D',\\
&c_2\notin P,\qquad P\cap U=\varnothing.
\end{aligned}}
\UebQuellen{\FormulaRefAuto{MogiljanskajaReserveMarkerFacts}[thm];
\FormulaRefAuto{MogiljanskajaReserveCoreFacts}[thm]}
\\ \addlinespace[.8ex]

\textbf{Zeilenparametrisierung und punktierte Verschiebung}
\UebFormel{\begin{aligned}
&\delta(n,j)\coloneqq b_{nj},\quad
\upsilon(n,j)\coloneqq b_{n+2,j},\\
&\theta\coloneqq\delta\circ\upsilon^{-1},\\
&H_n\coloneqq b_{n+1,0},\\
&\sigma\coloneqq(\SuccShift H)^{-1}.
\end{aligned}}
Die Nachfolgerverschiebung \(\SuccShift H\) wirkt längs der
injektiven Folge \(H\).
\UebQuellen{\FormulaRefAuto{MogiljanskajaRowParametrizationsDef}[def];
\FormulaRefAuto{MogiljanskajaThetaDef}[def];
\FormulaRefAuto{MogiljanskajaShiftSequenceDef}[def];
\FormulaRefAuto{MogiljanskajaSigmaDef}[def]}
&
\textit{Bijektionen und tatsächliche Bildmenge}
\UebFormel{\begin{aligned}
&\delta\colon\mathbb N^2\bij D,\qquad
\upsilon\colon\mathbb N^2\bij U,\\
&\theta\colon U\bij D,\qquad\sigma\colon U\bij V,\\
&H\colon\mathbb N\inj V,\qquad H_0=c_2,\\
&H[\mathbb N]=\{b_{n+1,0}\mid n\in\mathbb N\}
\subsetneq V.
\end{aligned}}
Hier ist \(\mathbb N^2=\mathbb N\times\mathbb N\).
\UebQuellen{\FormulaRefAuto{MogiljanskajaRowParametrizationsBijective}[thm];
\FormulaRefAuto{MogiljanskajaThetaBijective}[thm];
\FormulaRefAuto{MogiljanskajaSigmaBijective}[thm];
\FormulaRefAuto{MogiljanskajaShiftSequenceFacts}[thm];
\FormulaRefAuto{MogiljanskajaShiftSequenceImage}[thm]}
\\ \addlinespace[.8ex]

\textbf{Familien mit vorgegebener Bildmenge}
\UebFormel{\begin{aligned}
&\nu\coloneqq\sigma\circ\upsilon,\\
&I_{D'}\coloneqq(\{0\}\times\NatSeg2)
\cup(\{1\}\times\mathbb N^2),\\
&\omega((0,0))\coloneqq c_0,\qquad
\omega((0,1))\coloneqq c_1,\\
&\omega((1,(n,j)))\coloneqq\nu(n,j).
\end{aligned}}
\UebQuellen{\FormulaRefAuto{MogiljanskajaVParametrizationDef}[def];
\FormulaRefAuto{MogiljanskajaDPrimeIndexSetDef}[def];
\FormulaRefAuto{MogiljanskajaDPrimeParametrizationDef}[def]}
&
\textit{Die gewünschten Bildmengen werden exakt erreicht}
\UebFormel{\begin{aligned}
&\nu\colon\mathbb N^2\bij V,\qquad\nu[\mathbb N^2]=V,\\
&\omega\colon I_{D'}\bij D',\qquad\omega[I_{D'}]=D'.
\end{aligned}}
Zusammen mit \(a,b,\upsilon\) liegen damit ausdrücklich parametrisierte
Familien für \(L,D,U,V,D'\) vor.
\UebQuellen{\FormulaRefAuto{MogiljanskajaVParametrizationBijective}[thm];
\FormulaRefAuto{MogiljanskajaDPrimeParametrizationBijective}[thm];
\FormulaRefAuto{MogiljanskajaFamilyImages}[thm]}
\\
\end{BandUebersicht}
